% ============================================================
%  sections/02_related_work.tex
%  2. Công trình liên quan  (~0.5 trang)
% ============================================================

\section{Công Trình Liên Quan}
\label{sec:related}

% --- Phát hiện MGT tiếng Anh ---
\subsection*{Phát hiện MGT đơn ngôn ngữ}

Solaiman et al.~\cite{solaiman2019gpt2detector} giới thiệu bộ phát hiện
GPT-2 đầu tiên dựa trên RoBERTa, được fine-tuned trên văn bản do GPT-2
tạo ra. Dù đạt độ chính xác cao trên tiếng Anh, mô hình này không tổng
quát hóa sang ngôn ngữ khác do thiếu dữ liệu huấn luyện đa ngôn ngữ.
Mitchell et al.~\cite{mitchell2023detectgpt} đề xuất DetectGPT, một
phương pháp zero-shot dựa trên độ cong của hàm log-probability,
không cần nhãn huấn luyện nhưng đòi hỏi LLM ở mức trắng hộp (white-box).

% --- Phát hiện MGT đa ngôn ngữ ---
\subsection*{Phát hiện MGT đa ngôn ngữ}

Guo et al.~\cite{guo2023hcvariant} nghiên cứu phát hiện MGT trên tiếng
Trung, cho thấy các mô hình tiếng Anh thất bại hoàn toàn khi chuyển
sang ngữ cảnh MXH Trung Quốc (Weibo). Lý giải được đưa ra là đặc tính
bài đăng ngắn và ngữ nghĩa dày đặc ở cấp ký tự. Các nghiên cứu về
tiếng Việt~\cite{nguyen2023vsmec} và tiếng Ả Rập~\cite{alomari2022arabic}
cho thấy tương tự: hiện tượng ``Teencode'' trong tiếng Việt (viết tắt
có chủ đích, bỏ dấu) và cấu trúc hình thái phong phú trong tiếng Ả Rập
gây khó khăn đặc thù cho các detector.

% --- Phương pháp dựa trên Perplexity ---
\subsection*{Phát hiện dựa trên Perplexity}

Hans et al.~\cite{hans2024binoculars} và Equi et al.~\cite{equi2024ppl}
cho thấy văn bản do LLM tạo ra thường có perplexity (PPL) thấp hơn so
với văn bản con người khi được đánh giá bởi cùng một LLM, do phân phối
của chúng khớp với phân phối huấn luyện. Tuy nhiên, phương pháp này
nhạy cảm với nhiễu ngôn ngữ --- một đặc điểm phổ biến trong văn bản
MXH không chính thức.

% --- Vị trí bài báo này ---
Bài báo này bổ sung vào khoảng trống nghiên cứu bằng cách (i) xây dựng
bộ dữ liệu MXH thực tế cho cả bốn ngôn ngữ EN/VI/ZH/AR, (ii) đánh giá
đồng thời và có hệ thống ba họ phương pháp, và (iii) chứng minh tính khả
thi trên phần cứng GPU tiêu dùng (T4).
